\title{Group Sequence Policy Optimization}

\begin{abstract}
We present Group Sequence Policy Optimization (GSPO), a robust and effective reinforcement learning approach tailored for training large language models. Departing from prior methods that utilize token-level importance ratios, GSPO calculates importance ratios using entire sequence likelihoods and applies operations such as clipping, reward assignment, and optimization at the sequence level. Our results show that GSPO outperforms the GRPO algorithm in both training speed and performance, significantly enhances the stability of Mixture-of-Experts (MoE) RL training, and offers the promise of streamlining RL system architectures. The advancements delivered by GSPO have played a pivotal role in driving the substantial progress achieved in recent Qwen3 models.
\end{abstract}

\section{Introduction}
\label{sec:intro}

Reinforcement learning (RL) has become a foundational framework for advancing the scalability of language models \citep{o1,dpsk-r1,qwq32b,qwen3}. By employing large-scale RL approaches, these models acquire the ability to address complex challenges—including those found in high-level mathematics and coding—by engaging in more extended and nuanced reasoning.

Achieving effective RL scaling through increased computational resources fundamentally depends on ensuring consistent and reliable training stability. Yet, leading RL methods such as GRPO \citep{grpo} suffer from pronounced stability challenges when applied to training very large language models, frequently culminating in drastic and irrecoverable model failure \citep{qwen3, minimax-m1}. Such instability presents a significant barrier to advancing the performance limits of language models via extended RL training.

This study demonstrates that the primary cause of instability in GRPO lies in the incorrect use and consequent invalidation of importance sampling weights within its algorithmic framework. Such misapplication generates substantial high-variance noise during training, which not only accumulates as response sequences lengthen but is also exacerbated by the clipping process. Together, these factors drive the model toward collapse.

In response to these fundamental challenges, we introduce \textbf{Group Sequence Policy Optimization (GSPO)}, a novel reinforcement learning algorithm tailored for training large language models. The principal advancement of GSPO is its theoretically justified formulation of the importance ratio using sequence likelihood \citep{click}, consistent with the fundamental concept of importance sampling. Moreover, GSPO derives normalized rewards as the advantages among several responses to a given query, thereby preserving the consistency between rewards assigned at the sequence level and the optimization process.

Our empirical results indicate that GSPO substantially outperforms GRPO in terms of training stability, efficiency, and overall effectiveness. Notably, GSPO addresses the intrinsic stability issues commonly encountered in the RL training of large Mixture-of-Experts (MoE) models, obviating the necessity for intricate stabilization techniques and paving the way for streamlined RL systems. These advantages have played a pivotal role in driving the notable performance gains seen in the recent Qwen3 models. Looking forward, we anticipate that GSPO will serve as a durable and scalable foundation, supporting the further progress of large-scale RL training for language models.

\section{Preliminaries}

\paragraph{Notation}
Throughout this work, we represent an autoregressive language model with parameters $\theta$ as a policy $\pi_\theta$.
The symbol $x$ refers to a query, and $\mathcal{D}$ denotes the set of all queries. For a given query $x$ and its corresponding response $y$, the conditional probability of $y$ according to policy $\pi_\theta$ is written as $\pi_\theta (y | x)=\prod_{t=1}^{|y|} \pi_\theta (y_t | x, y_{<t} )$, where $|y|$ signifies the length of $y$ measured in tokens. Each query-response pair $(x, y)$ is evaluated by a verifier $r$, which assigns a reward value $r(x, y)$ within the interval $[ 0, 1 ]$.

\paragraph{Proximal Policy Optimization (PPO)} Proximal Policy Optimization (PPO) \citep{ppo} utilizes data sampled from the preceding policy $\pi_{\theta_\text{old}}$ and regulates the policy update by enforcing proximity to the prior policy via a clipping strategy. More precisely, PPO formulates policy optimization using the following objective (excluding the KL regularization term, which is omitted here for simplicity since it is not central to this work):

\begin{align}
\mathcal{J}_\text{PPO}(\theta) = \mathbb{E}_{ x \sim \mathcal{D},\, y \sim \pi_{\theta_\text{old}}( \cdot | x) }
\left[ \frac{1}{|y|} \sum_{t=1}^{|y|} 
\min \left( w_{t}(\theta) \widehat{A}_{t},  \, \mathrm{clip} \left( w_{t}(\theta), 1 - {\varepsilon}, 1 + {\varepsilon}\right) \widehat{A}_{t} \right)
\right],
\end{align}

Here, the importance ratio for token $y_t$ is given by $
w_{t}(\theta) = \frac{ \pi_{\theta} (y_{t} | x, y_{<t}) }{ \pi_{\theta_\text{old}} (y_{t} | x,y_{<t})} $. The advantage $\widehat{A}_{t}$ associated with $y_t$ is obtained using a separate value function, and $\varepsilon$ specifies the allowable range for clipping these importance ratios.

In practical applications, PPO faces a fundamental issue due to its strong dependence on the value model. In particular, the value model typically matches the policy model in scale, leading to significant demands on both memory and computational resources. Moreover, the performance of the algorithm critically depends on how accurate its value estimation is. Achieving a trustworthy value model is already a difficult endeavor; extending its applicability to handle longer sequences and more sophisticated problems compounds this difficulty even further.

\paragraph{Group Relative Policy Optimization (GRPO)} GRPO \citep{grpo} eliminates reliance on the value model by determining the relative advantage among responses belonging to the same query group. In particular, GRPO seeks to maximize the following objective:

\begin{align}
\label{equ:grpo}
\mathcal{J}_\text{GRPO}(\theta) = \mathbb{E}_{ x \sim \mathcal{D},\, \{y_i\}_{i=1}^G \sim \pi_{\theta_\text{old}}( \cdot | x) }
\left[ \frac{1}{G} \sum_{i=1}^{G} \frac{1}{|y_i|} \sum_{t=1}^{|y_i|} 
\min \left( w_{i,t}(\theta) \widehat{A}_{i,t},  \, \mathrm{clip} \left( w_{i,t}(\theta), 1 - {\varepsilon}, 1 + {\varepsilon}\right) \widehat{A}_{i,t} \right)
\right],
\end{align}

Here, $G$ denotes the total count of responses produced for each query $x$ (representing the group size). The terms $w_{i,t}(\theta)$ and $\widehat{A}_{i,t}$ respectively correspond to the importance ratio and the advantage for the token $y_{i,t}$, defined as:

\begin{align}
    w_{i,t}(\theta)=\frac{ \pi_{\theta} (y_{i,t} | x, y_{i,<t}) }{ \pi_{\theta_\text{old}} (y_{i,t} | x,y_{i,<t})},\quad\ 
    \widehat{A}_{i,t} = \widehat{A}_{i} = \frac{r(x, y_i) - \mathrm{mean} \left( \{ r(x, y_i) \}_{i=1}^G \right) }{ \mathrm{std} \left( \{ r(x, y_i) \}_{i=1}^G \right) },
\end{align}

respectively, with each token in $y_i$ being assigned the identical advantage $\widehat{A}_{i}$.

\section{Motivation}
\label{sec:motivation}

As models continue to increase in size, exhibit greater sparsity (such as in Mixture-of-Experts architectures), and produce longer responses, larger rollout batch sizes are required to ensure efficient hardware usage during reinforcement learning. To enhance sample efficiency, it is common to divide the large rollout batch into several smaller mini-batches for performing gradient updates. However, this approach creates an off-policy learning scenario, since the responses $y$ are drawn from a previous policy $\pi_{\theta_\text{old}}$ instead of the current policy $\pi_\theta$ that is currently being trained. Consequently, this justifies the use of clipping mechanisms in both PPO and GRPO, as these are designed to exclude samples that deviate excessively from the current policy when estimating gradients.

Although techniques such as clipping seek to address the off-policy mismatch, we uncover a deeper problem inherent to GRPO: \textit{its objective lacks a well-posed formulation}. This flaw is especially problematic during the training of large-scale models for tasks requiring extended responses, where it can trigger severe model degradation. The root of this issue lies in the inappropriate utilization of importance sampling weights within the GRPO framework. Importance sampling is designed to compute the expectation of a function $f$ with respect to a target distribution $\pi_\text{tar}$ by appropriately re-weighting data sampled from a behavior distribution $\pi_\text{beh}$:

\begin{align}
\mathbb{E}_{z \sim \pi_\text{tar}} \left[ f(z) \right] = \mathbb{E}_{z \sim \pi_\text{beh}} \left[ \frac{ \pi_\text{tar} (z) }{ \pi_\text{beh} (z)} \, f(z) \right].
\end{align}

Importantly, this approach depends on aggregating results across a large number of samples ($N \gg 1$) drawn from the behavior distribution $\pi_\text{beh}$ so that the importance weight $\frac{ \pi_\text{tar} (z) }{ \pi_\text{beh} (z)}$ can adequately compensate for differences between the distributions.

By comparison, GRPO incorporates the importance weight $\frac{ \pi_{\theta} (y_{i,t} | x, y_{i,<t}) }{ \pi_{\theta_\text{old}} (y_{i,t} | x,y_{i,<t})}$ at every token timestep $t$. This weighting is computed using only a single sampled token $y_{i,t}$ per next-token probability distribution $\pi_{\theta_\text{old}}( \cdot | x, y_{i, <t})$, which prevents it from effectively correcting for distribution shifts as intended. Rather than fulfilling its purpose, this method adds significant variance to the gradient estimates during training; over extended sequences, such variance becomes more pronounced, especially due to the use of a clipping procedure. In practice, we have found that these issues can result in catastrophic model collapse—an outcome from which recovery is extremely difficult. Efforts to restart training after collapse, such as rolling back to earlier model states, meticulously adjusting hyperparameters (like the clip threshold), increasing sequence lengths, or altering the nature of RL queries, typically prove unsuccessful.

The observation outlined above highlights a central flaw within GRPO's framework. Specifically, the ineffectiveness of the token-level importance weighting underscores an essential concept: \textbf{the granularity of the optimization objective must correspond to that of the reward signal}. Given that rewards are assigned per sequence, implementing off-policy correction on individual tokens is inherently misaligned. This realization leads us to discard the token-based approach and instead focus on employing importance weighting and conducting optimization at the \textit{sequence level}.

\section{Algorithm}

\subsection{GSPO: Group Sequence Policy Optimization}

Although the token-level importance weight $\frac{ \pi_{\theta} (y_{i,t} | x, y_{i,<t}) }{ \pi_{\theta_\text{old}} (y_{i,t} | x,y_{i,<t})}$ presents challenges within GRPO, we find that the \textit{sequence-level} importance weight $\frac{ \pi_{\theta} (y | x) }{ \pi_{\theta_\text{old}} (y | x)}$ holds significant theoretical relevance in language generation tasks. Specifically, this weight quantifies the extent to which a response $y$—drawn from $\pi_{\theta_\text{old}} (\cdot | x)$—differs from the distribution under $\pi_{\theta} (\cdot | x)$. Such a measure not only corresponds closely with sequence-level rewards but also provides a valuable metric for the operation of the clipping procedure.

Building on this simple insight, we introduce the \textbf{Group Sequence Policy Optimization (GSPO)} algorithm. GSPO optimizes policies through the following sequence-level objective:

\begin{align}
\mathcal{J}_\text{GSPO} (\theta) =
\mathbb{E}_{ x \sim \mathcal{D},\, \{y_i\}_{i=1}^G \sim \pi_{\theta_\text{old}}( \cdot | x) }
\left[ 
\frac{1}{G} \sum_{i=1}^{G}
\min \left( s_{i}(\theta)  \widehat{A}_{i},  \, \mathrm{clip} \left( s_{i}(\theta), 1 - {\varepsilon}, 1 + {\varepsilon} \right) \widehat{A}_{i} \right) 
\right],
\label{equ:gspo}
\end{align}

in which we employ the group-oriented estimation of advantage:

\begin{align}
\widehat{A}_{i} = \frac{r(x, y_i) - \mathrm{mean} \left( \{ r(x, y_i) \}_{i=1}^G \right) }{ \mathrm{std} \left( \{ r(x, y_i) \}_{i=1}^G \right) },
\end{align}

and establish the importance ratio $s_{i}(\theta)$ utilizing sequence likelihood as described in \citet{click}:

\begin{align}
s_{i}(\theta) = \left( \frac{ \pi_{\theta} (y_i | x) }{ \pi_{\theta_\text{old}} (y_i | x)} \right)^{\frac{1}{|y_i|}}
=
\exp \left( \frac{1}{|y_i|} \sum_{t=1}^{|y_i|} \log \frac{ \pi_{\theta} (y_{i,t} | x, y_{i,<t}) }{ \pi_{\theta_\text{old}} (y_{i,t} | x,y_{i,<t})} \right).
\end{align}

Consequently, GSPO implements clipping at the level of entire responses rather than on individual tokens, thereby removing excessively ``off-policy'' samples from the gradient computation; this approach aligns with both sequence-level reward assignment and optimization objectives. It is important to emphasize that we utilize length normalization in $s_{i}(\theta)$, which serves to minimize variance and maintain $s_{i}(\theta)$ within a consistent numerical interval. Without such normalization, changes in the likelihood of only a few tokens could induce large swings in the sequence-level importance ratio, and responses of varying lengths would necessitate different clipping thresholds. Furthermore, we observe that the clipping thresholds in GSPO generally differ substantially—in terms of order of magnitude—from those used in prior algorithms like GRPO, owing to the varying definitions of the importance ratios.

\subsection{Gradient Analysis}
\label{subsec:gradient}

The gradient of the GSPO objective may be determined as follows (note that clipping is excluded here for simplicity):

\begin{align}
\nabla_{\theta} \mathcal{J}_\text{GSPO} (\theta)
=&\ 
\nabla_{\theta} \mathbb{E}_{ x \sim \mathcal{D},\, \{y_i\}_{i=1}^G \sim \pi_{\theta_\text{old}}( \cdot | x) }
\left[ 
\frac{1}{G} \sum_{i=1}^{G} s_{i}(\theta) \widehat{A}_{i}
\right] \\
= &\ 
\mathbb{E}_{ x \sim \mathcal{D},\, \{y_i\}_{i=1}^G \sim \pi_{\theta_\text{old}}( \cdot | x) }
\left[ 
\frac{1}{G} \sum_{i=1}^{G}
s_{i}(\theta) \widehat{A}_{i}
\cdot \nabla_{\theta} \log s_{i}(\theta)
\right] \\
=&\ 
\mathbb{E}_{ x \sim \mathcal{D},\, \{y_i\}_{i=1}^G \sim \pi_{\theta_\text{old}}( \cdot | x) }
\left[ 
\frac{1}{G} \sum_{i=1}^{G} 
\left( \frac{ \pi_{\theta} (y_i | x) }{ \pi_{\theta_\text{old}} (y_i | x)} \right)^{\frac{1}{|y_i|}} \widehat{A}_{i} 
\cdot \frac{1}{|y_i|} \sum_{t=1}^{|y_i|} \nabla_{\theta} \log \pi_{\theta} (y_{i,t} | x, y_{i,<t}) 
\right].
\label{equ:gspo_grad}
\end{align}

To facilitate comparison, the gradient corresponding to the GRPO objective can be written as (with the identification $\widehat{A}_{i,t} = \widehat{A}_{i}$):

\begin{align}
\nabla_{\theta} \mathcal{J}_\text{GRPO} (\theta) 
=&\ 
\nabla_{\theta} \mathbb{E}_{ x \sim \mathcal{D},\, \{y_i\}_{i=1}^G \sim \pi_{\theta_\text{old}}( \cdot | x) }
\left[ \frac{1}{G} \sum_{i=1}^{G} \frac{1}{|y_i|} \sum_{t=1}^{|y_i|} 
w_{i,t}(\theta) \widehat{A}_{i,t}
\right] \\
=&\
\mathbb{E}_{ x \sim \mathcal{D},\, \{y_i\}_{i=1}^G \sim \pi_{\theta_\text{old}}( \cdot | x) }
\left[ \frac{1}{G} \sum_{i=1}^{G} \widehat{A}_{i} 
\cdot \frac{1}{|y_i|} \sum_{t=1}^{|y_i|} 
\frac{ \pi_{\theta} (y_{i,t} | x, y_{i,<t}) }{ \pi_{\theta_\text{old}} (y_{i,t} | x,y_{i,<t})} 
\nabla_{\theta} \log \pi_{\theta} (y_{i,t} | x, y_{i,<t})  
\right].
\end{align}

As a result, the primary difference between GSPO and GRPO is rooted in \textit{the manner in which they assign weights to the gradients of token log likelihoods}. GRPO employs an ``importance weight'' for each token, given by $\frac{ \pi_{\theta} (y_{i,t} | x, y_{i,<t}) }{ \pi_{\theta_\text{old}} (y_{i,t} | x,y_{i,<t})}$, which varies per token. These varying weights, which can fall within the intervals $(0, 1+\varepsilon]$ when $\widehat{A}_{i} >0$ or $[1-\varepsilon, +\infty)$ when $\widehat{A}_{i} < 0$, are not always close to unity and may cumulatively introduce instability during training. By contrast, GSPO assigns uniform weighting across all tokens in a given response, thereby avoiding the instability arising from the non-uniform weights in GRPO.

\subsection{GSPO-token: A Token-level Objective Variant}

In contexts such as multi-turn RL, it can be beneficial to perform advantage adjustment at a granularity finer than the sequence level. For this purpose, we present a token-level formulation of GSPO, referred to as \textbf{GSPO-token}, which enables the tailoring of advantages on a per-token basis:

\begin{align}
\mathcal{J}_\text{GSPO-token}(\theta)
=
\mathbb{E}_{ x \sim \mathcal{D},\, \{y_i\}_{i=1}^G \sim \pi_{\theta_\text{old}}( \cdot | x) }
\left[ 
\frac{1}{G} \sum_{i=1}^{G} \frac{1}{|y_i|} \sum_{t=1}^{|y_i|} 
\min \left( s_{i,t}(\theta) \widehat{A}_{i,t},  \, \mathrm{clip} \left( s_{i,t}(\theta), 1 - {\varepsilon}, 1 + {\varepsilon} \right) \widehat{A}_{i,t} \right) 
\right],
\label{equ:gspo-token}
\end{align}

in which

\begin{align}
s_{i,t}(\theta) 
= \mathrm{sg} \left[ s_{i}(\theta) \right]  \cdot \frac{ \pi_{\theta} (y_{i,t} | x, y_{i,<t}) }{ \mathrm{sg} \left[ \pi_{\theta} (y_{i,t} | x, y_{i,<t}) \right] },
\end{align}

and $\mathrm{sg}[\cdot]$ indicates extracting only the numerical value while blocking gradient propagation, analogous to the \texttt{detach} function in PyTorch. The gradient for GSPO-token is given by:

\begin{align}
\nabla_{\theta} \mathcal{J}_\text{GSPO-token}(\theta)
=&\ 
\nabla_{\theta} \mathbb{E}_{ x \sim \mathcal{D},\, \{y_i\}_{i=1}^G \sim \pi_{\theta_\text{old}}( \cdot | x) }
\left[ 
\frac{1}{G} \sum_{i=1}^{G}
\frac{1}{|y_i|} \sum_{t=1}^{|y_i|} 
s_{i,t}(\theta) \widehat{A}_{i,t}
\right] \\
=&\ 
\mathbb{E}_{ x \sim \mathcal{D},\, \{y_i\}_{i=1}^G \sim \pi_{\theta_\text{old}}( \cdot | x) }
\left[ 
\frac{1}{G} \sum_{i=1}^{G} s_i (\theta) \cdot \frac{1}{|y_i|} \sum_{t=1}^{|y_i|} 
\widehat{A}_{i,t} \frac{ \nabla_{\theta} \pi_{\theta} (y_{i,t} | x, y_{i,<t}) }{ \pi_{\theta} (y_{i,t} | x, y_{i,<t}) }
\right] \\
=&\ 
\mathbb{E}_{ x \sim \mathcal{D},\, \{y_i\}_{i=1}^G \sim \pi_{\theta_\text{old}}( \cdot | x) }
\left[ 
\frac{1}{G} \sum_{i=1}^{G}  \left( \frac{ \pi_{\theta} (y_i | x) }{ \pi_{\theta_\text{old}} (y_i | x)} \right)^{\frac{1}{|y_i|}}
\cdot \frac{1}{|y_i|} \sum_{t=1}^{|y_i|}
\widehat{A}_{i,t} \nabla_{\theta} \log \pi_{\theta} (y_{i,t} | x, y_{i,<t})  
\right].
\label{equ:gspo-token_grad}
\end{align}

Observe that the expression $\frac{ \pi_{\theta} (y_{i,t} | x, y_{i,<t}) }{ \mathrm{sg} \left[ \pi_{\theta} (y_{i,t} | x, y_{i,<t}) \right] }$ evaluates to 1, resulting in $s_{i,t}(\theta)$ being numerically the same as $s_{i}(\theta)$. By examining Equation~\eqref{equ:gspo} alongside~\eqref{equ:gspo-token}, as well as Equation~\eqref{equ:gspo_grad} versus~\eqref{equ:gspo-token_grad}, one can see that GSPO-token and GSPO coincide numerically in terms of their optimization objectives, clipping rules, and theoretical gradients when the token-level advantages for all tokens in a response $y_i$ are identical (i.e., $\widehat{A}_{i,t} = \widehat{A}_{i}$). Nevertheless, GSPO-token provides additional flexibility by allowing for token-specific adjustment of the advantages.

\section{Experiments and Discussion}

\subsection{Empirical Results}

We conduct experiments utilizing a cold-start model that has been fine-tuned from Qwen3-30B-A3B-Base. We present both the curves of training rewards and the performance metrics of the model on the AIME'24 (mean Pass@1 over 32 samples), LiveCodeBench (202410-202502, mean Pass@1 across 8 samples), and CodeForces (Elo Rating) evaluation suites. Throughout the RL training phase, each rollout batch is subdivided into four mini-batches to perform gradient updates. For GSPO, the clipping intervals in Equation~\eqref{equ:gspo} are chosen as 3e-4 (left) and 4e-4 (right). As a baseline, we benchmark against GRPO, where the clipping boundaries in Equation~\eqref{equ:grpo} are set to 0.2 (left) and 0.27 (right); these parameters have been carefully optimized to ensure comparability. It is important to highlight that GRPO requires the Routing Replay training procedure to achieve reliable MoE RL convergence, a subject explored further in \S~\ref{subsec:routing_replay}, whereas \textit{GSPO eliminates the reliance on this mechanism}.

Figure~\ref{fig:results} illustrates that training with GSPO maintains stable progress over time. Notably, we find that \textbf{GSPO enables ongoing improvements in performance by increasing training compute, systematically refreshing the query set, and lengthening the generated sequences}. In addition, compared to GRPO, GSPO exhibits greater training efficiency, achieving higher accuracy and better benchmark results using equivalent amounts of training compute and query consumption. Lastly, we have effectively utilized GSPO for reinforcement learning training on the latest Qwen3 models, providing compelling evidence for GSPO's effectiveness in harnessing the benefits of reinforcement learning scaling for large language models.

\subsection{Curious Observation on Clipping Fractions}

One fundamental difference between GSPO and GRPO lies in GSPO’s strategy of clipping at the sequence (response) level instead of the token level. As illustrated in Figure~\ref{fig:clipping}, this leads to a disparity of approximately two orders of magnitude in the proportions of clipped tokens when comparing GSPO with GRPO—a difference that persists even when the clipping ranges are varied. Interestingly, even though GSPO clips a much larger number of tokens and thus uses a smaller subset for training or gradient estimation, it nonetheless demonstrates superior training efficiency relative to GRPO. This somewhat surprising outcome—that greater clipping at the token level correlates with better efficiency—suggests that GRPO’s token-wise gradient estimates suffer from intrinsic noisiness and lack of efficiency for exploiting samples. Conversely, GSPO’s method, which leverages sequence-level gradients, yields a more dependable and robust learning signal.

\subsection{Benefit of GSPO for MoE Training}
\label{subsec:routing_replay}

\paragraph{Background}
Unlike standard RL training for dense architectures, MoE models, characterized by their sparse activation patterns, present distinct challenges concerning stability. Specifically, we observed that utilizing the GRPO algorithm with MoE models leads to pronounced \textit{expert-activation volatility}, impeding the proper convergence of RL training. In detail, after performing one or more gradient updates, the set of experts chosen to generate a particular response may shift substantially. For instance, considering the 48-layer Qwen3-30B-A3B-Base model, each RL gradient update leads to approximately 10\% of the experts, selected under the updated policy $\pi_{\theta}$, differing from those chosen by the preceding policy $\pi_{\theta_\text{old}}$ for the same rollout sample. This effect intensifies in deeper MoE architectures and results in significant oscillations of the token-level importance ratios $w_{i,t}(\theta) = \frac{ \pi_{\theta} (y_{i,t} | x, y_{i,<t}) }{ \pi_{\theta_\text{old}} (y_{i,t} | x,y_{i,<t})}$, undermining their reliability as detailed in \S~\ref{sec:motivation} and~\ref{subsec:gradient}. These instabilities consequently disrupt the standard convergence process in RL training.

\paragraph{Our Previous Approach} In addressing this issue, our earlier work utilized the \textbf{Routing Replay} training technique. In particular, we store the expert selection pattern (activated experts) determined by $\pi_{\theta_\text{old}}$ and, during the computation of the importance weights $w_{i,t}(\theta) = \frac{ \pi_{\theta} (y_{i,t} | x, y_{i,<t}) }{ \pi_{\theta_\text{old}} (y_{i,t} | x,y_{i,<t})}$, apply the same routing scheme to $\pi_{\theta}$. By doing so, the evaluations of $\pi_{\theta} (y_{i,t} | x, y_{i,<t})$ and $\pi_{\theta_\text{old}} (y_{i,t} | x,y_{i,<t})$ for a given token $y_{i,t}$ utilize identical expert networks. This approach stabilizes the token-level importance weights and guarantees that gradient-based updates act consistently on the same activated network. As shown in Figure~\ref{fig:routing_replay}, Routing Replay proves to be a vital component for achieving stable and reliable convergence during GRPO training of MoE models.

\paragraph{Benefit of GSPO} While Routing Replay allows GRPO to train MoE models with proper convergence, it introduces significant memory and communication costs by necessitating the reuse of routing patterns and may inadvertently restrict the model’s true capacity. By contrast, as illustrated in Figure~\ref{fig:results}, GSPO completely removes the reliance on Routing Replay, allowing for the standard computation of importance ratios $s_i(\theta)$, and achieves stable optimization and normal convergence behavior. The central idea underpinning GSPO is its emphasis on sequence-level likelihood (i.e., $\pi_{\theta} (y_i | x)$), rendering it less sensitive to fluctuations in the probability of individual tokens (i.e., $\pi_{\theta} (y_{i,t} | x, y_{i,<t})$). Since MoE models reliably preserve their language modeling strengths, the sequence-level likelihood exhibits minimal variation. In essence, GSPO directly addresses and mitigates the volatility in expert activation observed in MoE models, eliminating the need for auxiliary mechanisms like Routing Replay. This not only streamlines and fortifies the training pipeline, but also empowers the model to utilize its entire representational capacity unimpeded by artificial limitations.

\subsection{Benefit of GSPO for RL Infrastructure}

Due to differences in numerical precision between training engines (such as Megatron) and inference engines (like SGLang and vLLM), it is common in practice to recompute the likelihoods of generated responses under the previous policy $\pi_{\theta_\text{old}}$ using the training engine. In contrast, GSPO relies on sequence-level likelihoods for its optimization objective, rather than token-level likelihoods. Sequence-level metrics tend to be less sensitive to precision inconsistencies. As a result, GSPO enables the direct utilization of likelihoods computed by the inference engine, eliminating the necessity to recalculate them with the training engine. This feature is particularly advantageous in contexts such as partial rollout, multi-turn RL, or frameworks where training and inference are decoupled.

\section{Conclusion}

We introduce Group Sequence Policy Optimization (GSPO), a novel reinforcement learning method tailored for training large language models. Adhering to the core ideas of importance sampling, GSPO constructs importance ratios based on the probability of entire sequences and applies sequence-level clipping and rewards during the optimization process. Compared to GRPO, GSPO achieves greater training stability, efficiency, and overall performance, and proves particularly effective for scaling RL training in MoE architectures. This approach has underpinned significant advances observed in the most recent Qwen3 models. With GSPO serving as a robust, scalable foundation, we plan to further expand reinforcement learning approaches, anticipating major breakthroughs in artificial intelligence.

\end{document}